% TOP_PROBLEM: {"id":30007603,"problem_number":"LOCAL-30007603","title":"Liquidity regulation and solvency","category_id":21,"category":{"id":21,"name":"economics","display_name":"Economics","description":"Open economic theory statements, published research agendas, and proposed scoped specifications; question provenance and historical priority are qualified per record.","slug":"economics","order_index":21,"created_at":"2026-10-08T00:00:00Z"},"status":"research_question","external_url":null,"legacy_ids":[],"published":false,"statement":"In the explicitly specified setting: How should liquidity and capital rules interact when interest-rate losses, deposit flight and asset fire sales reinforce each other? The mathematical statement above fixes the target, assumptions and scope of this catalogue formulation.","economics":{"original_id":92,"field":"Banking, credit and financial stability","kind":"design","formalization_basis":"adapted_research_question","source_status":"research_agenda","priority_status":"earliest_found","priority_note":"This is the earliest explicit statement verified in this audit; first priority for the full catalogue formulation is not established.","earliest_verified_reference":{"authors":["Bank of England"],"title":"Bank of England Agenda for Research, 2025–2028","year":2026,"url":"https://www.bankofengland.co.uk/research/bank-of-england-agenda-for-research","doi":null,"locator":"Prudential Architecture / Regulatory framework; Resilience, paragraphs 1–2","evidence_summary":"Requests regulatory interaction and liquidity feedback analysis."},"current_reference":{"authors":["Bank of England"],"title":"Bank of England Agenda for Research, 2025–2028","year":2026,"url":"https://www.bankofengland.co.uk/research/bank-of-england-agenda-for-research","doi":null,"locator":"Prudential Architecture / Regulatory framework; Resilience, paragraphs 1–2","evidence_summary":"Requests regulatory interaction and liquidity feedback analysis."},"formalization_note":"Adapts an explicitly verified research-agenda passage. Mathematical objects, interventions, objectives and scope are proposed here; existing model-specific solutions are not relabeled unsolved."},"statusQualification":"Published question or research agenda verified at the cited locator. The mathematical specification is adapted for this catalogue and includes additional assumptions; source publication does not establish first-ever priority or certify this exact model remains unsolved. Adapts an explicitly verified research-agenda passage. Mathematical objects, interventions, objectives and scope are proposed here; existing model-specific solutions are not relabeled unsolved.","statusReviewedAt":"2026-10-08"}
% FIRST_POSTED: 2026-10-08
% ENTRY_KIND: statement_only
% !TeX program = lualatex
\documentclass[11pt]{article}
\usepackage{fontspec}
\usepackage[a4paper,margin=25mm]{geometry}
\usepackage{amsmath,amssymb,mathtools}
\usepackage{xurl}
\usepackage[hidelinks]{hyperref}
\newcommand{\sref}[2]{\hyperlink{source:#1:#2}{[S#2]}}
\setlength{\emergencystretch}{3em}
\begin{document}
% problemId:30007603
\section*{Liquidity regulation and solvency}\label{30007603}
\subsection{Definitions and mathematical statement}
Fix banks holding assets with specified durations, deposit liabilities and borrowing facilities. Let a joint rule \(a=(k,l)\) contain an equity requirement \(k\) and a liquid-asset requirement \(l\). A bank is solvent when its assets valued under the specified continuation allocation exceed liabilities; it is liquid when it can honor contemporaneous withdrawals without prohibited default. Interest-rate shocks change asset values, withdrawals affect financing needs and fire sales change market prices endogenously. Let \(W(a)\) be discounted social welfare in the resulting equilibrium, including intermediation services and crisis externalities. The task is to characterize the feasible joint optimum and its interaction term,
\[a^*\in\arg\max_{a\in A}W(a),\qquad\chi=\frac{\partial^2W(k,l)}{\partial k\,\partial l},\]
where \(A\) is the legally and operationally feasible rule set and derivatives are interpreted only away from binding-regime boundaries. Identify whether capital and liquidity requirements are substitutes or complements in each state, and quantify changes when bank duration and deposit composition differ. Specify deposit-insurance and lender-of-last-resort arrangements because they affect both constraints. Resolve equilibrium feedbacks rather than imposing fixed fire-sale prices. Where the objective is non-smooth, report finite policy changes and an identified set of welfare-maximizing rules. A capital loss and a liquidity shortfall must not be treated as interchangeable outcomes.

\par\textbf{Formulation and scope.} Adapts an explicitly verified research-agenda passage. Mathematical objects, interventions, objectives and scope are proposed here; existing model-specific solutions are not relabeled unsolved.

\par\textbf{Open-question provenance.} An explicit question or research agenda was verified at the source locator. This does not by itself certify that every later version remains unresolved \sref{30007603}{1}.

\par\textbf{Historical priority.} This is the earliest explicit statement verified in this audit; first priority for the full catalogue formulation is not established.

\subsection{Short English statement}
In the explicitly specified setting: How should liquidity and capital rules interact when interest-rate losses, deposit flight and asset fire sales reinforce each other? The mathematical statement above fixes the target, assumptions and scope of this catalogue formulation.

\subsection{Sources}
\begin{itemize}\raggedright
\item[\textnormal{[S1]}]\hypertarget{source:30007603:1}{} Bank of England. 2026. Bank of England Agenda for Research, 2025–2028. Prudential Architecture / Regulatory framework; Resilience, paragraphs 1–2.
\url{https://www.bankofengland.co.uk/research/bank-of-england-agenda-for-research}.
{\small\textit{Earliest verified question/agenda reference; historical priority qualified below. Requests regulatory interaction and liquidity feedback analysis.}}
\end{itemize}
\end{document}
